\chapter{Introducci\'on}
\label{sec:intro}

Con el objetivo de tener todos los ingredientes necesarios para la abstracci\'on
del juego de polic\'ias y ladrones, antes de adentrarnos a las reglas que rigen
una de estas partidas debemos sentar las bases que nos ayudar\'an a describir el
funcionamiento de este. En este caso, la noci\'on de un \textit{tablero} se
puede abstraer mediante una gr\'afica, dot\'ando al juego de toda la estructura
matem\'atica que esta conlleva. As\'i pues, las definiciones de las siguientes
tres secciones vienen de \cite{bondy} y \cite{diestel}.

\section{Gr\'aficas}

Dado un conjunto finito no vac\'io $V$, definimos una \indice{gr\'afica} $G$
como un par ordenado $(V_G,E_G)$ con $V_G$ un conjunto no vac\'io y $E_G$ un
subconjunto de $\binom{V}{1} \cup \binom{V}{2}$, donde
\[
  \binom{V}{k} = \l{e\subseteq V \colon\ \abs{e} = k}.
\]
Llamamos \textbf{v\'ertices}\index{v\'ertice} de $G$ a los elementos de $V_G$ y
\indice{orden} de $G$ a la cardinalidad de dicho conjunto. De forma similar,
llamamos \textbf{aristas}\index{arista} de $G$ y \indice{tama\~no} de $G$ a los
elementos y cardinalidad de $E_G$ respectivamente. Por convenci\'on, si
$e={x,y}$ con $x,y\in V_G$ distintos, denotamos por $xy$ a la arista $e$.
Decimos que una gr\'afica $G$ es \indiceSub{gr\'afica}{trivial} si tiene orden
$1$ y tama\~no $0$, de lo contrario, esta es \indiceSub{gr\'afica}{no trivial}.

En busca de facilitar el estudio de las propiedades de las gr\'aficas,
representamos estas estructuras a trav\'es de un dibujo en el plano, donde el
conjunto de v\'ertices corresponde a puntos distintos y cada arista es una curva
que une los puntos correspondientes a los v\'ertices en los que incide (ver
\cref{fig:grafica-dibujo}).

\begin{figure}[ht!]
  \centering
  \begin{tikzpicture}
    \begin{scope}[scale=1]
      \node (0) [vertex2, label=180:$v_1$] at (0,0){};
      \node (1) [vertex2, label=90:$v_2$]  at (2,0){};
      \begin{scope}[xshift=4cm]
        \node (2) [vertex2,label=90:$v_3$, below right=2cm of 1] {};
        \node (3) [vertex2,label=0:$v_4$ , above right=2cm of 2] {};
        \node (4) [vertex2,label=90:$v_5$, above right=2cm of 1] {};
      \end{scope}
      
      \draw [myloop,in=60,out=120,looseness=12]
            (0) to node[above]{$e_1$} (0);
      \draw [myloop,in=240,out=300,looseness=12]
            (2) to node[below]{$e_7$} (2);
      
      \draw [edge]  (0) to node [midway, above] {$e_2$} (1);
      \draw [edge]  (1) to node [midway, below left] {$e_3$} (2);
      \draw [edge]  (2) to node [midway, below right] {$e_4$} (3);
      \draw [edge]  (3) to node [midway, above right] {$e_5$} (4);
      \draw [edge]  (4) to node [midway, above left] {$e_6$} (1);
    \end{scope}
    \end{tikzpicture}
  \captionsetup{font=footnotesize}
  \caption{El dibujo de una gr\'afica $G$.}
  \label{fig:grafica-dibujo}
  \end{figure}

Dada una arista $e$ en $E_G$, decimos que un v\'ertice $u$ \textbf{incide}
\index{v\'ertice!incidente en arista} en $e$, y que $e$ incide en $u$, si $u\in
e$. Decimos que $e,f\in E_G$ distintas son \textbf{incidentes}
\index{arista!incidentes} si existe $u\in V_G$ incidente en ambas $e$ y $f$, y
decimos que estas son no incidentes en caso contrario. Dos v\'ertices distintos
$u$ y $v$ que inciden en una misma arista son \indiceSub{v\'ertice}{adyacentes},
de manera que $u$ es \textit{vecino} de $v$ y viceversa. En
\cref{fig:grafica-dibujo}, $v_1$ es incidente en $e_2$ y $e_2$ es incidente en
$v_1$, mientras que $e_1$ y $e_2$ son aristas incidentes y $v_1$ y $v_2$ son
v\'ertices adyacentes.

Si $e\in E_G$ es tal que existe $u\in V_G$ con $e=\l{u}$, decimos que $e$ es un
\indice{lazo} y $u$ tiene un lazo, mientras que $G$ es
\indiceSub{gr\'afica}{reflexiva} si $\l{u}\in E_G$ para cada $u\in V_G$.
Llamamos \indice{vecindad} de $v$ al conjunto $N_G(v)$ de v\'ertices de $G$
adyacentes a $v$; sus elementos ser\'an llamados vecinos de $v$. Adem\'as,
definimos la vecindad \indiceSub{vecindad}{cerrada} de $v$, $N_G[v]$, como la
uni\'on de $N_G(v)$ y $\l{v}$. Si $v$ no tiene lazo, denotamos por $d_G(v)$ a la
cardinalidad de $N_G(v)$, de lo contrario, $d_G(v) = \abs{N_G(v)} + 2$. Este
valor recibe el nombre de \indice{grado} de $v$. As\'imismo, representamos por
$\delta_G$ al grado m\'inimo sobre los v\'ertices de $G$, mientras que
$\Delta_G$ es el grado m\'aximo sobres los v\'ertices de $G$. De forma similar,
para $S\subseteq X$, denotamos por $N_G(S)$ a la vecindad del conjunto $S$,
definida como
\[
  N_G(S) = \bigcup_{x\in S} N_G(x).
\]
As\'i, en \cref{fig:grafica-dibujo}, $v_1$ tiene un lazo ($e_1$) y su grado es
$3$, donde $N_G(v_1)=\l{v_2}$. Adem\'as, $G$ no es reflexiva pues $v_2$ no tiene
un lazo. Si $S=\l{1,2}$, entonces $N_G(S)=\l{v_1,v_2,v_3,v_5}$. Adem\'as,
pasando sobre todos los v\'ertices de $G$ se obtiene que $\delta_G=2$ y
$\Delta_G=4$.

Dada $G$ una gr\'afica, decimos que $H$ es una \indice{subgr\'afica} de $G$ si
$V_H\subseteq V_G$ y $E_H\subseteq E_G$ donde para cada $e\in E_H$, $e\subseteq
V_H$. Adicionalmente, decimos que $H$ es una subgr\'afica
\indiceSub{subgr\'afica}{inducida} de $G$ si para cada $e\in E_G$ tal que
$e\subseteq V_H$, $e\in E_H$. Alternativamente, si $\varnothing\neq S\subseteq
V_G$, denotamos por $G[S]$ a la subgr\'afica inducida de $G$ cuyo conjunto de
v\'ertices es $S$ y cuyo conjunto de aristas tiene a todas las aristas
incidentes \'unicamente en v\'ertice en $S$. De forma similar, para $\varnothing
\ne E \subseteq E_G$, denotamos por $G[E]$ a la subgr\'afica de $G$ cuyo
conjunto de aristas es $E$ y su conjunto de v\'ertices son todos aquellos
v\'ertices tales que inciden en al menos una arista de $E$.

\begin{figure}[ht!]
  \centering
  \begin{tikzpicture}
    \begin{scope}[scale=1]
      \node (0) [vertex2, label=180:$v_1$,fill=red] at (0,0){};
      \node (1) [vertex2, label=90:$v_2$, fill=red]  at (2,0){};
      \begin{scope}[xshift=4cm]
      \node (2) [vertex2,label=90:$v_3$, fill=blue, below right=2cm of 1] {};
      \node (3) [vertex2,label=0:$v_4$, fill=blue, above right=2cm of 2] {};
      \node (4) [vertex2,label=90:$v_5$, above right=2cm of 1] {};
      \end{scope}
      
      \draw [myloop,in=60,out=120,looseness=12]
            (0) to node[above]{$e_1$} (0);
      \draw [myloop,in=240,out=300,looseness=12, blue]
            (2) to node[below]{$e_7$} (2);
      
      \draw [edge, dashed, red]  (0) to node [midway, above] {$e_2$} (1);
      \draw [edge]  (1) to node [midway, below left] {$e_3$} (2);
      \draw [edge, blue, dotted]  (2) to node [midway, below right] {$e_4$} (3);
      \draw [edge]  (3) to node [midway, above right] {$e_5$} (4);
      \draw [edge]  (4) to node [midway, above left] {$e_6$} (1);
    \end{scope}
    \end{tikzpicture}
  \captionsetup{font=footnotesize}
  \caption{La gr\'afica determinada por los v\'ertices y arista rojas, $H$, es
   una subgr\'afica de $G$. No obstante, esta no es inducida pues
   $G[\l{v_1,v_2}]$ tambi\'en incluye el lazo $e_1$. Por otro lado
   $H=G[\l{e_2}]$. En tanto, la gr\'afica inducida por los v\'ertices y aristas
   en azul, $H'$, s\'i es una subgr\'afica inducida de $G$, a saber,
   $H'=G[\l{v_3,v_4}]$.}
  \label{fig:grafica-dibujo2}
  \end{figure}

Un \indice{homomorfismo} de una gr\'afica $G$ en una gr\'afica $H$ es una
funci\'on $f \colon V_G\to V_H$ tal que $f(u)f(v)\in E_H$ para toda $uv\in E_G$.
Si adicionalmente pedimos que $f$ sea biyectiva y que $f(u)f(v)\in E_H$ si y
solo si $uv\in E_G$, decimos que $f$ es un \indice{isomorfismo} de gr\'aficas y
que $G$ y $H$ son \indiceSub{gr\'afica}{isomorfas}. Escribimos $G\cong H$ para
referirnos a que $G$ y $H$ son isomorfas. En \cref{fig:grafica-dibujo2}, existe
un homomorfismo de $G[\l{v_3,v_4}]$ en la gr\'afica $G[\l{v_3}]$, a saber, aquel
que fija $v_3$ y que manda $v_4$ a $v_3$. No obstante, estas gr\'aficas no son
isomorfas pues existe una \'unica funci\'on de $\l{v_3,v_4}$ en $\l{v_3}$, la
cual no es biyectiva. Por otro lado, $G[\l{v_1,v_2}] \cong G[\l{v_3,v_4}]$.

Decimos que $G$ es \indiceSub{gr\'afica}{completa} si $E_G=\binom{V}{2}$.
N\'otese que cualesquiera dos gr\'aficas completas de orden $n$ son isomorfas,
por lo que denotamos por $K_n$ a la gr\'afica completa de orden $n$. Una
gr\'afica $G$ es \indiceSub{gr\'afica}{bipartita} si $V_G$ puede partirse en dos
subconjuntos ajenos $X$ y $Y$ tales que toda arista incide en un v\'ertice en
$X$ y en un v\'ertice en $Y$. Tal partici\'on $(X,Y)$ es llamada una
bipartici\'on de la gr\'afica y $X$ y $Y$ son sus partes. Denotamos a una
gr\'afica bipartita $G$ con una partici\'on $(X,Y)$ por $G[X,Y]$. Si $G[X,Y]$ es
tal que ambos $X$ y $Y$ son distintos del vac\'io y todo v\'ertice de $X$ es
adyacente a todo v\'ertice de $Y$, decimos que $G$ es una gr\'afica
\indiceSub{gr\'afica}{bipartita completa}. En \cref{fig:grafica-dibujo2},
$X=\l{v_5}$ y $Y=\l{v_2,v_4}$ forman una partici\'on $(X,Y)$ de la gr\'afica
$G[\l{v_2,v_4,v_5}]$, que de hecho es bipartita completa. Por otra parte,
$X=\l{v_2}$ y $Y=\l{v_4}$ forman una partici\'on $(X,Y)$ de $G[\l{v_2,v_4}]$,
pero esta no es bipartita completa.

Una \indice{trayectoria} es una gr\'afica $G$ cuyos v\'ertices admiten un orden
total tal que dos v\'ertices son adyacentes en $G$ si y solo si son consecutivos
en el orden. Si $P$ es una trayectoria de orden $n$, suponemos sin p\'erdida de
generalidad que $V_P=\l{v_1,\dots,v_n}$, de manera que denotamos a las
trayectorias a trav\'es del orden de sus v\'ertices, de modo que
$P=\paren{v_1,\dots,v_n}$. De lo anterior es claro que, si $P$ y $P'$ son dos
trayectorias de orden $n$ con $n\ge 1$, entonces $P\cong P'$, donde el
isomorfismo es el que manda un v\'ertice de $P$ en la $i$-\'esima posici\'on al
v\'ertice de $P'$ en la $i$-\'esima posici\'on. As\'i pues, denotamos por $P_n$
a la trayectoria de orden $n$. Similarmente, un \indice{ciclo} consiste de una
trayectoria en la que su primer y \'ultimo v\'ertice del orden son adyacentes.
Por el mismo argumento, existe un \'unico ciclo de orden $n$, al que denotamos
$C_n$. La \indice{longitud} de una trayectoria o un ciclo es el tama\~no de
dicha gr\'afica. Diremos que un ciclo o una trayectoria son pares o impares
acorde a la paridad de su longitud. Por otra parte, sean $G$ una gr\'afica y
$u,v\in V_G$. Decimos que $W=(v_1,v_2,\dots, v_n)$ es un $uv$-\indice{camino} en
$G$ si $v_iv_{i+1}\in E_G$ para cada $i\in\seg{1}{n-1}$. Si adicionalmente
$v_i\neq v_j$ cada para $i\neq j$, decimos que $P=(v_1,\dots, v_n)$ es una
$uv$-trayectoria en $G$. Si $i\leq j$, denotamos por $v_iPv_j$ al segmento de la
trayectoria $P$ que comienza en $v_i$ y termina en $v_j$. Adicionalmente, si
$v_jx\in E_G$, denotamos por $v_iPv_jx$ a la subgr\'afica de $G$ inducida por
las aristas del segmento $v_iPv_j$ y la arista $v_jx$. Analogamente se define un
ciclo en $G$ como $C=(v_1,\dots,v_n,v_1)$, donde $v_{i}v_{i+1}\in E_G$ para cada
$i\in\seg{1}{n-1}$, $v_nv_1\in E_G$ y $v_i\neq v_j$ para cada $i\neq j$.
Continuando con los ejemplos en \cref{fig:grafica-dibujo2}, $P=(v_2,v_5,v_4)$ es
una $v_2v_4$-trayectoria en $G$, y es tal que $P\cong P_3$. Por otro lado,
$v_2Pv_4v_3v_2$ es un ciclo par de longitud $4$ isomorfo a $C_4$.

Pasamos ahora a demostrar el siguiente teorema b\'asico que nos da una
condici\'on suficiente para que una gr\'afica contenga a un ciclo como
sugr\'afica.

\begin{teorema}
\label{teo:grado-ciclo}
    Si $G$ es una gr\'afica tal que $\delta_G \ge 2$, entonces $G$ contiene un
    ciclo.
\end{teorema}
  
\begin{proof}
  Si $G$ tiene un lazo en un v\'ertice $v$, $C=(v,v)$ es un ciclo en $G$.
  Sup\'ongase entonces que $G$ no tiene lazos. Sea $P=(v_1,\dots,v_n)$ una
  trayectoria de longitud m\'axima en $G$. Por ser $P$ de longitud m\'axima,
  $N_G(v_n)\subseteq \seg{v_1}{v_{n-1}}$. Adem\'as, $v_nv_{n-1}\in E_G$ y no hay
  lazo en $v_n$, por lo que
  \[
      2\leq \delta_G \leq d_G(v_n) = \abs{N_G(v_n)} = 1 + 
      \abs{N_G(v_n)\cap \seg{v_1}{v_{n-2}}}.
  \]
  As\'i, existe $i\in\seg{1}{n-2}$ tal que $v_nv_i\in E_G$. Luego,
  $v_iPv_nv_i$ es un ciclo en $G$.

\end{proof}

El teorema anterior nos permite demostrar un resultado similar que nos asegura
la existencia de un ciclo en $G$. Dadas dos gr\'aficas $G$ y $H$, con conjuntos
de v\'ertices no necesariamente ajenos, definimos la \indice{uni\'on} de $G$ y
$H$, $G \cup H$, como la gr\'afica cuyo conjunto de v\'ertices es $V_G\cup V_H$
y con conjunto de aristas $E_G\cup E_H$.

\begin{proposicion}
\label{pro:trayectoria-ciclo}
    Sean $P=(v_1,\dots,v_n)$ y $Q=(u_1,\dots,u_m)$ trayectorias distintas en $G$
    tales que $v_1=u_1$ y $v_n=u_m$. Entonces $P\cup Q$ contiene un ciclo.
\end{proposicion}
  
\begin{proof}
  Como $E_P\neq E_Q$ por ser $P$ y $Q$ trayectorias distintas, se puede
  considerar $H=G[E_P\triangle E_Q]$. Se afirma que $\delta_H\ge 2$. En efecto,
  sea $v\in V_H$. Si $v=v_1$ o $v=v_n$, dado que el grado de $v$ es $1$ tanto en
  $P$ como en $Q$, debe suceder que la arista de $P$ incidente en $v$ es
  distinta de la arista de $Q$ incidente en $v$, por lo que $d_H(v)\ge 2$.
  
  As\'i pues, sup\'ongase que $v$ es un v\'ertice distinto de los extremos de
  $P$. Como $v\in V_H$, existe, sin p\'erdida de generalidad, $e\in E_P$
  incidente en $v$ y $e\notin E_Q$. Al ser $v$ un v\'ertice distinto de los
  extremos de $P$, existe $h\neq e$ arista de $P$ incidente en $v$. Si $h\notin
  E_Q$, entones $d_H(v)\ge 2$. De lo contrario, nuevamente por ser $v$ distinto
  de los extremos de $Q$, existe $e'\in E_Q$ distinta de $h$ incidente en $v$.
  N\'otese que $e'\neq e$ pues $e\notin E_Q$. Luego, como $v$ tiene grado $2$ en
  $P$, $e'\notin E_P$, por lo que $e'\in E_H$. As\'i, $d_H(v)\ge 2$.

  Esto implica que $\delta_H\ge 2$, por lo que, en virtud
  d\cref{teo:grado-ciclo}, $P\cup Q$ contiene un ciclo.

\end{proof}

Pasamos ahora a definir una propiedad fundamental de las gr\'aficas. Una
gr\'afica $G$ es \indiceSub{gr\'afica}{conexa} si para cualquier par de
v\'ertices $u$ y $v$, existe una $uv$-trayectoria en $G$. Adem\'as, llamaremos
\indice{componente conexa} a una subgr\'afica inducida maximal de $G$ con la
propiedad de ser conexa. Si $u$ y $v$ son dos v\'ertices en una gr\'afica conexa
$G$, denotamos por $d_G(x,y)$ a la longitud de la $xy$-trayectoria m\'as
peque\~na en $G$. Decimos que $G$ es ac\'iclica si esta no tiene ciclo alguno
como una subgr\'afica. As\'i, en \cref{fig:grafica-dibujo} es claro que $G$ es
conexa, mientras que $G[\l{v_2,v_4}]$ no lo es. No obstante, $G$ no es aciclica,
mientras que $G[\l{v_2,v_4}]$ s\'i lo es. Adem\'as, algunas distancias entre
v\'ertices de $G$ son $d_G(v_1,v_1)=0$, $d_G(v_1,v_2)=1$ y $d_G(v_2,v_4)=2$.

El siguiente teorema caracteriza las componentes conexas de una gr\'afica $G$
tal que $\Delta_G\leq 2$.

\begin{teorema}
\label{teo:compconexa}
  Si $G$ es una gr\'afica con $\Delta_G\leq 2$ y $C$ es una componente conexa de
  $G$, entonces $C$ es una trayectoria o un ciclo.
\end{teorema}

\begin{proof}
  Sea $P=(v_1,\dots,v_n)$ una trayectoria de longitud m\'axima en $C$. Si $1=n$,
  entonces $C$ es la gr\'afica trivial, y por tanto una trayectoria, o es la
  gr\'afica con un v\'ertice y un lazo, por lo que es un ciclo. Luego, si $1\neq
  n$, al ser $P$ de longitud m\'axima se tiene que $N_G(v_n)\subseteq
  \seg{v_1}{v_{n-1}}$. N\'otese que para cada $i\in\seg{2}{n-2}$ y para
  cualquier v\'ertice $u$ distinto de $v_{i-1}$ y $v_{i+1}$, $uv_i\notin E_G$
  pues, de lo contrario,
  \[
      2\ge \Delta_G \ge d_G(v_i) \ge 3,  
  \]
  lo cual es una contradicci\'on. As\'i, $N_G(v_n)\subseteq \l{v_{n-1},v_1}$ y
  $N_G(v_1)\subseteq \l{v_n,v_2}$, mientras que $N_G(v_i)=\l{v_{i-1},v_{i+1}}$
  para cualquier $i\in\seg{2}{n-2}$. Adem\'as, ning\'un v\'ertice en $C$ tiene
  lazos pues $\Delta_G\leq 2$, por lo que $C$ es una trayectoria si
  $v_1v_n\notin E_G$, o es un ciclo si $v_1v_n\in E_G$.

\end{proof}

\section{\'Arboles}

Una gr\'afica $G$ es un \textbf{\'arbol}\index{arbol@\'arbol} si esta es conexa
y ac\'iclica (ver \cref{fig:grafica-dibujo3}). Por ser ac\'iclica, autom\'aticamente tenemos que $G$ no tiene
lazos. M\'as a\'un, de \cref{pro:trayectoria-ciclo} se tiene el siguiente
corolario inmediato.

\begin{corolario}
\label{cor:unitray}
  Sea $T$ un \'arbol, y sean $u$ y $v$ v\'ertices de $T$. Entonces existe una
  \'unica $uv$-trayectoria en $T$.
\end{corolario}

\Cref{teo:grado-ciclo} implica que en un \'arbol no trivial existe al menos un
v\'ertice de grado $1$. Llamamos \textbf{hojas}\index{hoja} a los v\'ertices de
grado $1$ en un \'arbol.

\begin{figure}[H]
  \centering
  \begin{tikzpicture}
    \begin{scope}[scale=1]
      \node (0) [vertex2, label=180:$v_1$, fill=red] at (0,0){};
      \node (1) [vertex2, label=90:$v_2$, fill=red]  at (2,0){};
      \begin{scope}[xshift=4cm]
      \node (2) [vertex2,label=90:$v_3$, fill=red, below right=2cm of 1] {};
      \node (3) [vertex2,label=0:$v_4$, fill=red, above right=2cm of 2] {};
      \node (4) [vertex2,label=90:$v_5$, fill=red, above right=2cm of 1] {};
      \end{scope}
      
      \draw [myloop,in=60,out=120,looseness=12]
            (0) to node[above]{$e_1$} (0);
      \draw [myloop,in=240,out=300,looseness=12]
            (2) to node[below]{$e_7$} (2);
      
      \draw [edge, red]  (0) to node [midway, above] {$e_2$} (1);
      \draw [edge, red]  (1) to node [midway, below left] {$e_3$} (2);
      \draw [edge, red]  (2) to node [midway, below right] {$e_4$} (3);
      \draw [edge, red]  (3) to node [midway, above right] {$e_5$} (4);
      \draw [edge]  (4) to node [midway, above left] {$e_6$} (1);
    \end{scope}
    \end{tikzpicture}
  \captionsetup{font=footnotesize}
  \caption{Las aristas de rojo inducen la gr\'afica $G[\l{e_2,e_3,e_4,e_5}]$, la
  cual es un \'arbol y tiene a $v_1$ y $v_5$ como hojas.}
  \label{fig:grafica-dibujo3}
  \end{figure}

Dada una gr\'afica $G$, podemos a\~nadir o eliminar v\'ertices a $G$ para
generar una nueva gr\'afica a partir de $G$. Si $u\in V_G$, y
$V_G\setminus\l{u}\neq \varnothing$ denotamos por $G-u$ a la subgr\'afica
inducida $G[V_G\setminus \l{u}]$. Similarmente, si $e$ es una arista con
$e=\l{x,y}$, decimos que $H$ se obtiene de $G$ al subdividir la arista $e$ si
$V_H=V_G\cup\l{v_e}$ con $v_e\notin V_G$ y $E_H=E_G\cup \l{v_ex, v_ey} \setminus
\l{xy}$.

Una observaci\'on inmediata es que si $x$ es una hoja de $T$ y $P$ es una
trayectoria en $T$ que no utiliza a $x$ entonces $P$ es una trayectoria en
$T-x$. As\'i, por \cref{cor:unitray}, la $uv$-trayectoria en $T$ con $u\neq
x\neq v$ no utiliza a $x$ pues $d_T(x)=1$ y el grado (en P) de todo v\'ertice en
$P$ distinto de los extremos es $2$. Adem\'as, dado que $T$ es ac\'iclica, es
claro que $T-x$ tambi\'en lo es, por lo que tiene lugar el siguiente corolario.

\begin{corolario}
\label{cor:eliminacion-arbol}
  Si $T$ es un \'arbol y $x$ es una de sus hojas, entonces $T-x$ es un \'arbol.
\end{corolario}

De la misma manera, si $e$ es una arista de $T$, la gr\'afica obtenida de la
subdivisi\'on de $e$, $T'$ tambi\'en da lugar a un \'arbol pues todas las
trayectorias en $T$ que no usaban a $e$ siguen siendo trayectorias en $T'$,
mientras que las que la usaban se convierten en trayectorias de un v\'ertice
m\'as al usar las dos aristas incidentes en $v_e$ (ver
\cref{fig:grafica-dibujo4}). En tanto, de haber un ciclo en $T'$, este
utilizar\'ia a $v_e$ y a sus dos aristas, por lo que un ciclo de un v\'ertice
menos existir\'ia en $T$, contradiciendo que este es un \'arbol.

\begin{figure}[ht!]
  \centering
  \begin{tikzpicture}
    \begin{scope}[scale=1]
      \node (0) [vertex2, label=180:$v_1$, fill=blue] at (0,0){};
      \node (1) [vertex2, label=90:$v_2$, fill=blue]  at (2,0){};
      \begin{scope}[xshift=4cm]
      \node (2) [vertex2,label=180:$v_3$, fill=red, below right=2cm of 1] {};
      \node (3) [vertex2,label=0:$v_4$, fill=red, above right=2cm of 2] {};
      \node (4) [vertex2,label=90:$v_5$, fill=blue!50!white, above right=2cm of 1] {};
      \node (5) [vertex2,label=315:$v_6$, fill=red, above right=0.75cm of 2] {};
      \end{scope}
      
      \draw [myloop,in=60,out=120,looseness=12]
            (0) to node[above]{$e_1$} (0);
      \draw [myloop,in=240,out=300,looseness=12]
            (2) to node[below]{$e_7$} (2);
      
      \draw [edge, blue]  (0) to node [midway, above] {$e_2$} (1);
      \draw [edge]  (1) to node [midway, below left] {$e_3$} (2);
      \draw [edge, red]  (2) to node [midway, above left] {$v_3v_6$} (5);
      \draw [edge, red]  (5) to node [midway, above left] {$v_6v_4$} (3);
      \draw [edge]  (3) to node [midway, above right] {$e_5$} (4);
      \draw [edge, blue!50!white, dashed]  (4) to node [midway, above left] {$e_6$} (1);
    \end{scope}
    \end{tikzpicture}
  \captionsetup{font=footnotesize}
  \caption{Al \'arbol $G[\l{e_2,e_6}]$ se le elimina la hoja $v_5$, resultado en
  un \'arbol. Por otra parte, al \'arbol $G[\l{e_4}]$ se le subdivide la arista
  $e_4$ para generar otro \'arbol.}
  \label{fig:grafica-dibujo4}
\end{figure}
As\'i pues, se llega al siguiente corolario.

\begin{corolario}
\label{cor:subdivision-arbol}
  Si $T$ es un \'arbol y $e$ es una de sus aristas, entonces la gr\'afica obtenida
  al subdividir $e$ es un \'arbol.
\end{corolario}

\section{Emparejamientos}

Dada una gr\'afica $G$, decimos que $S\subseteq V_G$ es un conjunto
\indice{dominante} de $G$ si todo v\'ertice fuera de $S$ es adyacente a al menos
un v\'ertice de $S$; al m\'inimo de las cardinalidades sobre todos los conjuntos
dominantes de $G$ le llamamos \indice{n\'umero de dominancia} y lo denotamos por
$\gamma(G)$. Similarmente, $M\subseteq E_G$ es un \indice{emparejamiento} en $G$
si $M$ no contiene lazos y cualesquiera dos aristas $e$ y $f$ en $M$ son no
incidentes. Si $M$ es un emparejamiento, decimos que los dos extremos de una
arista en $M$ est\'an \textit{emparejados} bajo $M$, mientras que, para
cualquier v\'ertice incidente en una arista de $M$, decimos que el v\'ertice
est\'a \indiceSub{v\'ertice}{saturado} por $M$. Un \textbf{emparejamiento
perfecto}\index{emparejamiento!perfecto} es un emparejamiento que satura a todos
los v\'ertices de la gr\'afica, mientras que un emparejamiento es
\indiceSub{emparejamiento}{m\'aximo} si satura la mayor cantidad de v\'ertices
de entre todos los emparejamientos. Denotamos por $\alpha'(G)$ al
\textbf{n\'umero de emparejamiento}\index{emparejamiento!n\'umero de} de $G$, es
decir, la cardinalidad de un emparejamiento m\'aximo en $G$.

Dado un emparejamiento $M$ en una gr\'afica $G$, definimos una
$M$-\textbf{trayectoria} \indiceSub{trayectoria}{alternante} en $G$ como una
trayectoria cuyas aristas alternan, siguiendo el orden de la trayectoria, entre
aristas en $M$ y en $E_G\setminus M$. Si tanto el inicio como el final de la
trayectoria no est\'an saturados por $M$, decimos que esta es una
$M$-\textbf{trayectoria} \indiceSub{trayectoria}{aumentable} (ver
\cref{fig:grafica-dibujo5}).

\begin{figure}[ht!]
  \centering
  \begin{tikzpicture}
    \begin{scope}[scale=1]
      \node (0) [vertex2, label=180:$v_1$, fill=blue] at (0,0){};
      \node (1) [vertex2, label=90:$v_2$, node split half=90]  at (2,0){};
      \begin{scope}[xshift=4cm]
      \node (2) [vertex2,label=90:$v_3$, fill=red, below right=2cm of 1] {};
      \node (3) [vertex2,label=0:$v_4$, fill=red, above right=2cm of 2] {};
      \node (4) [vertex2,label=90:$v_5$, fill=red, above right=2cm of 1] {};
      \end{scope}
      
      \draw [myloop,in=60,out=120,looseness=12]
            (0) to node[above]{$e_1$} (0);
      \draw [myloop,in=240,out=300,looseness=12]
            (2) to node[below]{$e_7$} (2);
      
      \draw [edge, blue]  (0) to node [midway, above] {$e_2$} (1);
      \draw [edge]  (1) to node [midway, below left] {$e_3$} (2);
      \draw [edge, red]  (2) to node [midway, below right] {$e_4$} (3);
      \draw [edge]  (3) to node [midway, above right] {$e_5$} (4);
      \draw [edge, red]  (4) to node [midway, above left] {$e_6$} (1);
    \end{scope}
    \end{tikzpicture}
  \captionsetup{font=footnotesize}
  \caption{El conjunto de aristas azules, $\l{e_2}$, es un emparejamiento $M$ en
  $G$ que empareja a $v_1$ y $v_2$ y los satura. M\'as a\'un, $M$ no es m\'aximo
  pues $M^*=\l{e_4,e_6}$ es un emparejamiento de cardinalidad mayor. Este s\'i
  es m\'aximo en $G$, m\'as no es un emparejamiento perfecto. No obstante, $M^*$
  s\'i es un emparejamiento p\'erfecto en $G-v_1$. Adem\'as, $(v_1,v_2,v_5)$ es
  una $M^*$-trayectoria alternante, mientras que $(v_3,v_4)$ es una
  $M$-trayectoria aumentable. Finalmente, $\l{v_1,v_2,v_3}$ es un conjunto
  dominante, pero $\gamma(G)=2$ pues $\l{v_2,v_4}$ tambi\'en lo es.}
  \label{fig:grafica-dibujo5}
\end{figure}

\newpage
El siguiente teorema relaciona la existencia de $M$-trayectorias aumentables con
los emparejamientos m\'aximos.

\begin{teorema}[Teorema de Berge]
\label{teo:berge}
  Un emparejamiento $M$ en una gr\'afica $G$ es m\'aximo si y solo si $G$ no
  contiene $M$-trayectorias aumentables.
\end{teorema}

\begin{proof}
  Sea $M$ un emparejamiento en $G$. Sup\'ongase que $G$ contiene una
  $M$-trayectoria aumentable, $P$. V\'ease que $M'=M\triangle E_P$ es un
  emparejamiento en $G$. En efecto, $M'$ no contiene lazos pues ni $M$ ni $P$
  los contienen. Por otro lado, para $e,f\in M'$, si ambas aristas se encuentran
  en $M\setminus E_P$, entonces estas son no incidentes. Si ambas aristas
  est\'an en $E_P\setminus M$, entonces estas son no incidentes pues $P$ alterna
  entre aristas en $M$ y en $E\setminus M$. Finalmente, si $e\in M\setminus E_P$
  y $f\in E_P\setminus M$, estas son no incidentes ya que, de lo contrario, $f$
  y $e$ inciden en un v\'ertice no saturado por $M$ (si dicho v\'ertice es el
  incial o final de $P$) o inciden en un v\'ertice saturado por $M$ a trav\'es
  de una arista en $P$ distinta de $e$, contradiciendo que $M$ es un
  emparejamiento.
  
  As\'i, $M'$ es un emparejamiento y, dado que $P$ es $M$-alternante, este tiene
  exactamente una arista m\'as en $E\setminus M$ de las que tiene en $M$, por lo
  que $\abs{M'} = \abs{M} + 1$, contradiciendo que $M$ era un emparejamiento
  m\'aximo en $G$.

  Rec\'iprocamente, sup\'ongase que $M$ es un emparejamiento en $G$ y que este
  no es m\'aximo. Sea $M^*$ un emparejamiento m\'aximo en $G$, de manera que
  $\abs{M^*} > \abs{M}$. Sea $H=G[M\triangle M^*]$.

  Obs\'ervese que cada v\'ertice de $H$ tiene grado (en $H$) $1$ o $2$ pues este
  puede ser incidente en a lo m\'as una arista de $M$ y a lo m\'as una arista de
  $M^*$. As\'i, $\Delta_H\leq 2$, por lo que, en virtud de
  \cref{teo:compconexa}, cada componente conexa de $H$ es un ciclo o es una
  trayectoria. M\'as a\'un, si es un ciclo, este es par pues alterna entre
  aristas en $M$ y $M^*$ por ser estos emparejamientos, mientras que la
  trayectoria alterna entre aristas en $M$ y $M^*$ por la misma raz\'on.

  Como $\abs{M^*}>\abs{M}$, la subgr\'afica $H$ contiene m\'as aristas en $M^*$
  que en $M$, por lo que debe existir en $H$ una componente conexa $P$ que
  consiste de una trayectoria que alterna entre aristas en $M^*$ y $M$ y cuyas
  primera y \'ultima arista est\'an en $M^*\setminus M$. As\'i, ambos el
  v\'ertice inicial y final de $P$ no est\'an saturados por $M$, de manera que
  $P$ es una $M$-trayectoria aumentable.

\end{proof}

En gr\'aficas bipartitas es sencillo dar una condici\'on necesaria y suficiente
para que exista un emparejamiento que sature a todos los v\'ertices en una de
las partes de la bipartici\'on.

\begin{teorema}[Teorema de Hall]
\label{teo:hall}
  Una gr\'afica bipartita $G[X,Y]$ tiene un emparejamiento que satura todos los
  v\'ertices en $X$ si y solo si
  \[
      \abs{N(S)}\ge \abs{S}
  \]
  para cada $S\subseteq X$.
\end{teorema}

\begin{proof}
  Sea $G[X,Y]$ una gr\'afica bipartita con un emparejamiento $M$ que satura
  todos los v\'ertices en $X$. Sea $S\subseteq X$. Los v\'ertices en $S$ est\'an
  emparejados bajo $M$ a distintos v\'ertices en $N(S)$, por lo que
  $\abs{N(S)}\ge \abs{S}$.

  Sup\'ongase ahora que $G[X,Y]$ es una gr\'afica bipartita y sup\'ongase que
  $G$ no tiene un emparejamiento que sature a todos los v\'ertices en $X$. Sea
  $M^*$ un emparejamiento m\'aximo en $G$ y $u\in X$ no saturado por $M^*$. Sea
  $Z$ el conjunto de todos los v\'ertices tales que, si $z\in Z$, existe una
  $M^*$-trayectoria alternante que comienza en $u$ y termina en $z$. Como $M^*$
  es un emparejamiento m\'aximo, se sigue d\cref{teo:berge} que $u$ es el
  \'unico v\'ertice en $Z$ no saturado por $M^*$. Sean $R=X\cap Z$ y $B=Y\cap
  Z$.

  \begin{figure}[ht!]
    \centering
    \begin{tikzpicture}
      \begin{scope}[scale=1]
        \node (0) [vertex2, label=180:$u$, fill=blue] at (0,0){};
        \node (1) [vertex2, label=90:$v_2$, fill=red]  at (2,0){};
        \begin{scope}[xshift=4cm]
        \node (2) [vertex2,label=90:$v_3$, fill=blue, below right=2cm of 1] {};
        \node (3) [vertex2,label=0:$v_4$, fill=red, above right=2cm of 2] {};
        \node (4) [vertex2,label=90:$z$, fill=blue, above right=2cm of 1] {};
        \end{scope}
        
        \draw [edge]  (0) to node [midway, above] {$e_1$} (1);
        \draw [edge]  (1) to node [midway, below left] {$e_2$} (2);
        \draw [edge, verdeAzulado]  (2) to node [midway, below right] {$e_3$} (3);
        \draw [edge]  (3) to node [midway, above right] {$e_4$} (4);
        \draw [edge, verdeAzulado]  (4) to node [midway, above left] {$e_5$} (1);
      \end{scope}
      \end{tikzpicture}
    \captionsetup{font=footnotesize}
    \caption{La gr\'afica $G$ es bipartita con la partici\'on $(X,Y)$ dada por
    la coloraci\'on de los v\'ertices, de manera que los v\'ertices azules
    est\'an en $X$ y los rojos en $Y$. Las aristas verdes forman un
    emparejamiento m\'aximo $M^*$ por \cref{teo:berge}. Adem\'as, utilizando las
    trayectorias $M^*$-alternantes se obtiene que $Z=\l{u,v_2,z,v_4,v_3}$.
    As\'i, $R=\l{u,z,v_3}$ y $B=\l{v_2,v_4}$.}
    \label{fig:grafica-dibujo6}
  \end{figure}

  Claramente los v\'ertices en $R\setminus\l{u}$ est\'an emparejados bajo $M^*$
  a v\'ertices en $B$. As\'i, $\abs{B} \ge \abs{R} - 1$ , mientras que cada
  v\'ertice en $B$ est\'a saturado por $M^*$ ya que, de lo contrario, la $M^*$
  trayectoria alternante de $u$ a $y$ ser\'ia $M^*$-aumentable, contradiciendo
  que $M^*$ es m\'aximo por \cref{teo:berge}. Luego, $\abs{B} = \abs{R} - 1$.
  Por otro lado, $B\subseteq N(R)$ pues, dado $y\in B$, el pen\'ultimo v\'ertice
  de una $uy$-trayectoria $M^*$-alternante est\'a en $R$. En tanto, todo
  v\'ertice $y$ en $N(R)$ tiene un vecino $x$ que es alcanzado por $u$ a
  trav\'es de una $M^*$-trayectoria alternante $P$ cuyo \'ultimo v\'ertice es
  $x$. Si $y$ es un v\'ertice interno de $P$, $y\in B$, mientras que si no,
  $uPxy$ es una trayectoria $M^*$-alternante pues $x$ est\'a saturado por $M^*$,
  de manera que $xy\notin M^*$. Por tanto, $B=N(R)$ y
  \[
    \abs{N(R)} = \abs{B} = \abs{R} - 1 < \abs{R},
  \]
  de manera que la condici\'on dada por el teorema falla para $S=R$.

\end{proof}

\section{Gr\'aficas de fichas}

Dada una gr\'afica $G$, estamos interesados en obtener una estructura que nos
permita estudiar \textit{configuraciones} de $G$. No es de sorprender que esa
estructura que buscamos es tambi\'en una gr\'afica. As\'i, en \cite{fabila},
Fabila et al. introducen el concepto de \textit{gr\'afica de fichas} de una
gr\'afica $G$ y estudian sus propiedades b\'asicas como gr\'afica en t\'erminos
de las propiedades de $G$. Por tanto, las siguientes definiciones y resultados
se obtienen de dicho trabajo.

Dado $X$ un conjunto y $k\leq \abs{X}$, decimos que $S\subseteq X$ es un
$k$-conjunto de $X$ si $\abs{S}=k$. Adem\'as, dados dos conjuntos $A$ y $B$,
recordemos que la \indice{diferecia sim\'etrica} de $A$ y $B$, denotada por
$A\triangle B$, se define como
\[
  A \triangle B = (A\setminus B) \cup (B\setminus A).
\]
As\'i, dada una gr\'afica $G$ sin lazos y un entero $k$ tal que $1\leq k\leq
\abs{V_G}$, definimos la \textbf{gr\'afica de k fichas}\index{gr\'afica!de
fichas} de $G$, $\t{k}{G}$, cuyo conjunto de v\'ertices consiste de todos los
$k$-conjuntos de $V_G$, mientras que dos de estos $k$-conjuntos, $A$ y $B$, son
adyacentes si y solo si $A\triangle B\in E_G$ (ver \cref{fig:grafica-fichas}).

\begin{figure}[H]
\centering
  \begin{tikzpicture}
    \triadatres{1.8cm}{(360/6)*0}
    \triadatres{1.8cm}{(360/6)*1}
    \triadatres{1.8cm}{(360/6)*2}
    \triadatres{1.8cm}{(360/6)*3}
    \triadatres{1.8cm}{(360/6)*4}
    \triadatres{1.8cm}{(360/6)*5}

    \foreach \i in {0,2,4}
      \node (\i) [bola2]  at ({(360/6)*\i}:1.8){};
    \foreach \i in {1,3,5}
      \node (\i) [bola2] at ({(360/6)*\i}:1.8){};
    
    \foreach \i in {0,...,5}
      \draw [edge] let \n1={int(mod(\i+1,6))} in (\i) to (\n1);

    \begin{scope}[shift={(2)}]
      \node (P1) [punto3,fill=black] at ({90}:0.4) {};
      \node (P2) [punto3,fill=black] at (0:0) {};
    \end{scope}

    \begin{scope}[shift={(3)}]
      \node (P3) [punto3,fill=black] at ({90}:0.4) {};
      \node (P4) [punto3,fill=black] at ({90 + 360/3}:0.4) {};
    \end{scope}

    \begin{scope}[shift={(4)}]
      \node (P5) [punto3,fill=black] at (0:0) {};
      \node (P6) [punto3,fill=black] at ({90 + 360/3}:0.4) {};
    \end{scope}

    \begin{scope}[shift={(5)}]
      \node (P7) [punto3,fill=black] at ({90 + (360/3)*2}:0.4) {};
      \node (P8) [punto3,fill=black] at ({90 + 360/3}:0.4) {};
    \end{scope}

    \begin{scope}[shift={(0)}]
      \node (P9) [punto3,fill=black] at ({90 + (360/3)*2}:0.4) {};
      \node (P10) [punto3,fill=black] at (0:0) {};
    \end{scope}

    \begin{scope}[shift={(1)}]
      \node (P11) [punto3,fill=black] at ({90 + (360/3)*2}:0.4) {};
      \node (P12) [punto3,fill=black] at ({90}:0.4) {};
    \end{scope}
  \end{tikzpicture}
  \captionsetup{font=footnotesize}
  \caption{Se muestra $\t{2}{G}$, la gr\'afica de $2$ fichas de la gr\'afica
  bipartita completa $G[X,Y]$ con $\abs{X}=1$ y $\abs{Y}=3$. Dentro de cada
  v\'ertice de $\t{2}{G}$ se dibuj\'o a $G$ junto con un $2$-conjunto de sus
  v\'ertices identificado por los puntos. N\'otese que los v\'ertices adyacentes
  a un v\'ertice de $\t{2}{G}$ corresponden a aquellas configuraciones que se
  alcanzan al ``deslizar'' una ficha a lo largo de una arista de $G$.}
  \label{fig:grafica-fichas}
\end{figure}

Una observaci\'on inmediata de la definici\'on es que, si $G$ es una gr\'afica
sin lazos, $G$ es isomorfa a su gr\'afica de $1$ ficha, es decir, $G\cong
\t{1}{G}$. Esto es inmediato al ver que los $1$-conjuntos de $V_G$ son los
unitarios de los v\'ertices y que dos $1$-conjuntos son adyacentes en $\t{1}{G}$
si y solo si su diferencia sim\'ertica es una arista en $G$, es decir, si estos
son adyacentes en $G$.

El siguiente lema reduce los casos a estudiar para las gr\'aficas de fichas.

\begin{lema}
  \label{lem:isomorfismo}
      Sean $k$ y $n$ enteros positivos. Si $G$ es una gr\'afica sin lazos de
      orden $n$ y $k < n$, entonces $\t{k}{G} \cong \t{n-k}{G}$.
  \end{lema}

  \begin{proof}
    Se propone la funci\'on $\varphi \colon V_{\t{k}{G}} \to V_{\t{n-k}{G}}$
    dada por la regla de correspondencia $\varphi(A) = V\setminus A$. Dicha
    funci\'on est\'a bien definida pues $\abs{V\setminus A} = n-k$. Por tanto,
    resta demostrar que dicha funci\'on da un isomorfismo entre $\t{k}{G}$ y
    $\t{n-k}{G}$.

    Si $AB\in E_{\t{k}{G}}$, entonces $A\triangle B =\l{a,b}\in E_G$ con $a\in
    A\setminus B$ y $b\in B\setminus A$. As\'i,
    \begin{align*}
      f(A)\triangle f(B) 
        &= \paren{V\setminus A}\triangle \paren{V\setminus B} \\
        &= \paren{V\setminus A \setminus \paren{V\setminus B}} \cup 
           \paren{V\setminus B \setminus \paren{V\setminus A}}\\
        &= \paren{\paren{V\setminus A} \cap B} \cup \paren{\paren{V\setminus B}  \cap A}\\
        &= \l{b} \cup \l{a} = \l{a,b}\in E_G.
    \end{align*}
    Luego, $f(A)f(B)\in E_{\t{n-k}{G}}$. Por la misma raz\'on, si $f(A)f(B)\in
    E_{\t{n-k}{G}}$, entonces $AB\in E_{\t{k}{G}}$, por lo que se concluye que
    $\t{k}{G} \cong \t{n-k}{G}$.

  \end{proof}

\section{Polic\'ias y ladrones}

Llegamos finalmente a la definici\'on del juego de polic\'ias y ladrones. Como
objeto de estudio matem\'atico, la versi\'on de polic\'ias y ladrones que hoy en
d\'ia conocemos es de alrededor de 1980. El par\'ametro que nos dedicaremos a
estudiar a lo largo del texto, el \textit{n\'umero policiaco}, fue introducido
por Aigner y Fromme en \cite{aignerDAM8}. Por otra parte, en \cite{bonato},
Bonato y Nowakowski realizan una revisi\'on documental de todos los resultados
que se tienen respecto al juego.

As\'i pues, \indice{polic\'ias y ladrones} es un juego de
evasi\'on-persecuci\'on sobre una gr\'afica en el cual dos jugadores, $C$ y $R$
(del ingl\'es, \textit{cops \& robbers}), toman turnos alternadamente para
conseguir su objetivo. Por una parte, $C$ controla a una cantidad fija $c$ de
polic\'ias tales que estos ocupan v\'ertices no necesariamente distintos de $G$
y, en cada turno, todos se mueven a trav\'es de una arista incidente al
v\'ertice en el que se encontraban previamente. Por otro lado, $R$ controla a un
\'unico ladr\'on cuyas reglas de movimiento son las mismas que las descritas
para los polic\'ias de $C$. Por tanto, el objetivo de $C$ es que al menos uno de
sus polic\'ias ocupe el mismo v\'ertice que el ladr\'on, momento en el que
diremos que dicho polic\'ia \indice{captura} a $R$ y termina el juego de
inmediato con la victoria de $C$. En tanto, el objetivo de $R$ es evitar su
captura, de manera que $R$ gana el juego si puede hacer esto de forma
indefinida. Por convenci\'on, $C$ comienza el juego al escoger posiciones
iniciales para cada uno de sus polic\'ias, mientras que, posteriormente, $R$ en
su primer turno decide en qu\'e v\'ertice colocar a su ladr\'on.

Diremos que un jugador tiene una \indice{estrategia ganadora} si cuenta con un
conjunto de reglas a seguir tal que, sin importar lo que haga su contrincante,
este es capaz de ganar el juego. Por tanto, $C$ tiene una estrategia ganadora si
puede eventualmente capturar a $R$ tras un n\'umero finito de movimientos y sin
importar los movimientos que este realice, mientras que $R$ tiene una estrategia
ganadora si puede evadir su captura sin importar los movimientos que realice
$C$.

La versi\'on antes descrita es la llamada versi\'on ``activa'' del juego.   Hay
una versi\'on ``pasiva'' en la que, en su turno, cada polic\'ia y el ladr\'on
pueden pasar, de modo que estos mantienen su posici\'on. No obstante, esta
variante del juego puede emularse con la versi\'on activa al jugar sobre
gr\'aficas reflexivas, de manera que, si alg\'un polic\'ia o el ladr\'on buscan
pasar, estos se mueven a trav\'es del lazo asociado al v\'ertice en el que se
encuentran. As\'i, a lo largo del trabajo examinaremos el juego tradicional
\'unicamente sobre gr\'aficas reflexivas a\'un cuando nos referiremos a dichas
gr\'aficas como sus versiones sin lazos.

Dada una gr\'afica $G$ de orden $n$, es claro que si $C$ tiene $n$ polic\'ias,
entonces tiene una estrategia ganadora, a saber, en su primer turno coloca cada
polic\'ia en un v\'ertice distinto de $G$, de manera que, sin importar cu\'al
v\'ertice escoja $R$, este pasa a ser capturado inmediatamente. Por tanto, cabe
preguntarse, dada una gr\'afica $G$, cu\'al es el m\'inimo n\'umero de
polic\'ias que $C$ necesita para capturar a $R$. A este par\'ametro lo
llamaremos el \indice{n\'umero policiaco} de $G$ y lo denotaremos por $c(G)$.

Dadas una gr\'afica $G$ y $H$ una de sus subgr\'aficas inducidas, decimos que
$H$ es un \indice{retracto} de $G$ si existe un homomorfismo $f$ de $G$ en $H$
tal que $f(x)=x$ para cualquier $x\in V_H$. A trav\'es de esta noci\'on, podemos
dar una cota inferior para el n\'umero policiaco de una gr\'afica en funci\'on
de sus retractos.

\newpage

\begin{lema}
\label{lem:retracto-cota}
  Si $H$ es un retracto de $G$, entonces $c(H)\leq c(G)$.  
\end{lema}

\begin{proof}
  Sup\'ongase que $k$ polic\'ias tienen una estrategia ganadora en $G$ y sea $f$
  un homomorfismo de $G$ a $H$ tal que $f(x)=x$ para cada $x\in V_H$. Se
  consideran dos partidas indepdendientes en paralelo de polic\'ias y ladrones:
  una jugada sobre $G$ y otra sobre $H$. El juego en $H$ puede considerarse como
  un juego sobre $G$ pues $H$ es una subgr\'afica inducida de $G$. La estrategia
  ganadora de $C$ con $k$ polic\'ias en $G$ puede no ser suficiente para
  capturar a $R$ en $H$, por ejemplo, puede que $R$ tenga que salir de $H$ para
  ser capturado en la parte de $G$ que no est\'a en $H$.

  As\'i pues, se considera la siguiente estrategia de sombras: los polic\'ias
  jugando la partida en $G$ siguen la estrategia usual, mientras que los
  polic\'ias jugando la partida en $H$ realizan los movimientos dados por las
  im\'agenes bajo $f$ de los polic\'ias en $G$. Por simplicidad, se denota por
  $f(c)$ a las imagenes de las posiciones de los polic\'ias en $G$. As\'i, si un
  polic\'ia $c_i$ en $G$ se mueve de un v\'ertice $u$ a $v$, entonces un
  polic\'ia $f(c_i)$ en $H$ se mueve de $f(u)$ a $f(v)$, el cual es un
  movimiento v\'alido por ser $f$ un homomorfismo.

  Se afirma que la estrategia descrita es ganadora para los polic\'ias en $H$.
  En efecto, consid\'erese que los polic\'ias juegan en $G$ y $R$ est\'a
  restringido a $H$. Si los polic\'ias est\'an a punto de ganar en $G$, debe
  suceder que el v\'ertice en el que se encuentra $R$ y todos sus vecinos, tanto
  en $H$ como en $G$, son adyacentes a un polic\'ia. Luego, si $u$ es el
  v\'ertice en el que se encuentra $R$ y $v$ es un v\'ertice adyacente a este,
  el homomorfismo $f$ manda la arista $uv\in E_G$ a la arista $uf(v)\in E_H$.
  Luego, como $f$ es homomorfismo, si $c_i$ es el polic\'ia en un v\'ertice
  adyacente a $v$ en $G$, $f(c_i)$ es un polic\'ia adyacente a $f(v)$ en $H$. Lo
  mismo sucede con el polic\'ia en $G$ que est\'a en un v\'ertice adyacente al
  v\'ertice en el que se encuentra $R$. Por tanto, todos los posibles
  movimientos de $R$ son tales que un polic\'ia en $H$ le captura en a lo m\'as
  su siguiente turno, de modo que $k$ polic\'ias bastan para capturar a $R$ en
  $H$.

\end{proof}

\begin{lema}
\label{lem:cotasup-dominante}
  Si $G$ es una gr\'afica, entonces $c(G)\leq \gamma(G)$.
\end{lema}

\begin{proof}
  Sea $S\subseteq V_G$ un conjunto dominante de $G$ de cardinalidad $\gamma(G)$.
  As\'i, si $C$ tiene $\gamma(G)$ polic\'ias, este pone uno de ellos en cada
  v\'ertice de $S$ en su primer turno. Sup\'ongase que $R$ se posiciona en $v$.
  Si $v\in S$, hay un polic\'ia en $v$, por lo que el juego termina. De lo
  contrario, existe una arista que incide en $v$ y es tal que su otro extremo,
  $u$ est\'a en $S$ por ser este dominante. As\'i, el polic\'ia que se encuentra
  en $u$ procede a capturar a $R$ en su siguiente turno.

\end{proof}

A manera de exploraci\'on, demostramos el siguiente teorema para \'arboles.

\begin{teorema}
\label{teo:copnum-arbol}
  Si $T$ es un \'arbol, entonces $c(T)=1$.
\end{teorema}

\begin{proof}
  Si $T$ es trivial, el resultado es claro. Por tanto, sup\'ongase que es no
  trivial. Se dar\'a una estrategia ganadora para $C$ que controla a un \'unico
  polic\'ia. Se denotar\'a por $C$ al polic\'ia que controla $C$, de modo que
  este se posicione en un v\'ertice arbitrario de $T$. De la misma forma, se
  denotar\'a por $R$ al ladr\'on que controla $R$. En los siguientes turnos,
  dado que $T$ es un \'arbol, por \cref{cor:unitray} se tiene que existe una
  \'unica trayectoria en $T$ del v\'ertice en el que se encuentra $C$ al
  v\'ertice en el que se encuentra $R$. As\'i, una estrategia para $C$ es mover
  al segundo v\'ertice de la \'unica trayectoria entre su posici\'on actual y la
  posici\'on de $R$.

  Esto implica que, posterior al turno de $R$, la distancia entre $C$ y $R$ no
  aumenta nunca. M\'as a\'un, dicha distancia debe disminuir eventualmente pues
  si $R$ evita que eso pase, este llega eventualmente a una hoja, que existe por
  ser $T$ no trivial. Luego, inductivamente se tiene que la distancia entre $C$
  y $R$ eventualmente disminuye de $1$ a $0$, por lo que $C$ captura a $R$.

\end{proof}

Dada una gr\'afica $G$ sin lazos y $k\ge 1$, podemos considerar un juego de
polic\'ias y ladrones sobre su gr\'afica de $k$ fichas, recordando que el juego
se desarrolla sobre la versi\'on reflexiva de $\t{k}{G}$. No obstante, en lugar
de trabajar sobre $\t{k}{G}$ para obtener estrategias ganadoras, podemos
recurrir a la misma gr\'afica para plantear una modificaci\'on del juego
tradicional, de manera que se puede emular una partida cl\'asica sobre
$\t{k}{G}$. Est\'a variante la introduce Fern\'andez-Vel\'azquez en
\cite{fernandez}. Para una partida tradicional de polic\'ias y ladrones sobre
$\t{k}{G}$, consid\'erese una partida de polic\'ias y ladrones por equipos en
$G$, donde cada polic\'ia de $C$ y el ladr\'on tienen un equipo de $k$ fichas
tales que, por cada polic\'ia y el ladr\'on, sus $k$ fichas determinan un
$k$-conjunto de $V_G$. Adem\'as, en su turno, cada polic\'ia y el ladr\'on
escogen una de sus $k$ fichas y la mueven de un v\'ertice a otro a trav\'es de
una arista, de manera que posterior a su movimiento estos siguen determinando un
$k$-conjunto de $V_G$. Por otra parte, la acci\'on de que un polic\'ia o el
ladr\'on pase el turno representa, en la gr\'afica de fichas, que el respectivo
polic\'ia o el ladr\'on no mueva alguna de sus fichas en $G$. En esta ocasi\'on,
$C$ gana si en su turno logra que las $k$ fichas de uno de sus polic\'ias
determinen el mismo $k$-conjunto que las $k$ fichas del ladr\'on, mientras que
el ladr\'on gana si puede evitar indefinidamente que esto suceda. En algunas
ocasiones nos referiremos a los v\'ertices de $\t{k}{G}$ como
\textbf{configuraciones}\index{configuraci\'on}.

De este punto en adelante, denotamos por $c_i$ al $i$-\'esimo polic\'ia de $C$,
mientras que al ladr\'on, controlado por $R$, lo denotamos por $r$. Nos interesa
distinguir entre las fichas de $r$, por lo que, para cada $1\leq j\leq k$
denotamos por $r_j$ a la $j$-\'esima ficha de $r$. Asimismo, distinguimos las
fichas de cada polic\'ia $c_i$, de modo que, para cada $1\leq j\leq k$,
denotamos por $c_{i,j}$ a la $j$-\'esima ficha del $i$-\'esimo polic\'ia.
Decimos que un polic\'ia $c_i$ \indice{domina} una configuraci\'on $S\in
V_{\t{k}{G}}$ si este se encuentra en $S$ o en una configuraci\'on adyacente a
$S$. Adem\'as, decimos que una ficha $c_{i,j}$ de $c_i$
\textbf{captura}\index{captura!de fichas} a $r_j$ si estas se encuentran en el
mismo v\'ertice. En algunos casos abusaremos del vocabulario al decir que $c_i$
\textit{captura} a $r_j$ si una ficha de $c_i$ le captura. Finalmente, $c_i$
\textbf{captura en la gr\'afica de fichas}\index{captura!en la gr\'afica de
fichas} a $r$ si ambos determinan el mismo $k$-conjunto de $V_{G}$. Nuevamente
abusando del vocabulario, diremos que $C$ \textit{captura} a $R$ si $c_i$
captura a $r$ para alg\'un \'indice $i$.